% Généré par tools/build.py — ne pas éditer à la main
\section{Nombres réels et intervalles}

\fichehead{jeu.~19~nov.~2026}{Ensembles de nombres~: $\N,\Z,\D,\Q,\R$}
\begin{fiche}
\qitem{Échauffement · Flash}{NUM}{Donner l'écriture décimale de $\dfrac38$ et une valeur approchée de $\dfrac56$ à $10^{-3}$ près.}{$0{,}375$\,; $0{,}833$.}{}
\qitem{Cœur · Méthode}{NUM}{Donner le plus petit ensemble ($\N,\Z,\D,\Q,\R$) contenant chacun de ces nombres~: $-4$\,; $0{,}25$\,; $\dfrac13$\,; $\sqrt2$\,; $\dfrac63$.}{$-4\in\Z$\,; $0{,}25\in\D$\,; $\dfrac13\in\Q$\,; $\sqrt2\in\R$\,; $\dfrac63=2\in\N$.}{Le plus petit ensemble est la clé~: $\dfrac63$ est un entier naturel.}
\qitem{Cœur + · Vrai/faux}{NUM}{Vrai ou faux\,? Justifier. Le nombre $\dfrac13$ est un nombre décimal.}{Faux. S'il existait $a$ et $n$ entiers avec $\dfrac13=\dfrac{a}{10^n}$, alors $10^n=3a$ serait divisible par $3$~: impossible (la somme des chiffres de $10^n$ vaut $1$).}{}
\qitem{Retour · Flash}{SUB}{Calculer $(x+1)(x-3)$ pour $x=2$.}{$-3$.}{}
\qitem{Fil rouge · Flash}{SOL}{Combien de faces, sommets et arêtes possède un prisme droit à base triangulaire\,?}{$5$ faces ($2$ triangles, $3$ rectangles), $6$ sommets, $9$ arêtes.}{}
\end{fiche}

\fichehead{lun.~23~nov.~2026}{Encadrements décimaux et arrondis}
\begin{fiche}
\qitem{Échauffement · Flash}{NUM}{Arrondir $3{,}14159$ à $10^{-2}$ près puis à $10^{-3}$ près.}{$3{,}14$ puis $3{,}142$.}{Le chiffre suivant décide~: $1<5$ puis $5\geqslant5$.}
\qitem{Cœur · Méthode}{NUM}{Encadrer $\sqrt2$ par deux décimaux d'amplitude $10^{-2}$.}{$1{,}41^2=1{,}9881<2<2{,}0164=1{,}42^2$~: $1{,}41<\sqrt2<1{,}42$.}{}
\qitem{Cœur + · Lecture de code}{ALG}{Quelle est la valeur finale de \code{k}\,? Que peut-on en déduire sur $\sqrt2$\,?\par\smallskip{\ttfamily\small \hspace*{0.0em}k = 10\\ \hspace*{0.0em}while (k/10)**2 \textless{} 2:\\ \hspace*{2.4em}k = k + 1\par}}{$k=15$~: $1{,}4^2=1{,}96<2$ et $1{,}5^2=2{,}25\geqslant2$, donc $1{,}4<\sqrt2\leqslant1{,}5$.}{}
\qitem{Retour · Flash}{QP}{$5$ kg de riz coûtent $12\,500$ Ar. Combien coûtent $8$ kg\,?}{$\dfrac{12\,500}{5}\times8=20\,000$ Ar.}{}
\qitem{Fil rouge · Vrai/faux}{DEM}{Vrai ou faux\,? Justifier. Un quadrilatère dont les diagonales ont même longueur est un rectangle.}{Faux. Contre-exemple~: un trapèze isocèle a des diagonales de même longueur sans être un rectangle.}{}
\end{fiche}

\fichehead{mar.~24~nov.~2026}{Nombres irrationnels~: $\frac13$ n'est pas décimal, $\sqrt2$ est irrationnel}
\begin{fiche}
\qitem{Échauffement · Flash}{CAL}{Écrire $\sqrt{12}$ et $\sqrt{50}$ sous la forme $a\sqrt b$ avec $b$ entier le plus petit possible.}{$\sqrt{12}=2\sqrt3$\,; $\sqrt{50}=5\sqrt2$.}{}
\qitem{Cœur · Tri}{NUM}{Parmi $0{,}333\ldots$\,; $\sqrt9$\,; $\sqrt2$\,; $\pi$\,; $\dfrac{22}{7}$, lesquels sont rationnels\,?}{$0{,}333\ldots=\dfrac13$\,; $\sqrt9=3$\,; $\dfrac{22}7$~: rationnels. $\sqrt2$ et $\pi$ sont irrationnels.}{$\dfrac{22}{7}$ n'est qu'une valeur approchée de $\pi$, pas égale à $\pi$.}
\qitem{Cœur + · Raisonnement}{DEM}{On suppose $\sqrt2=\dfrac pq$ avec $\dfrac pq$ irréductible, d'où $p^2=2q^2$. Que dire de la parité de $p$, puis de $q$\,? Quelle contradiction\,?}{$p^2$ pair, donc $p$ pair, $p=2k$\,; alors $4k^2=2q^2$, $q^2=2k^2$~: $q$ pair. Contradiction avec la fraction irréductible.}{}
\qitem{Retour · Tableau}{STA}{Valeurs $1\,;\,2\,;\,3\,;\,4$ d'effectifs $5\,;\,8\,;\,4\,;\,3$. Donner l'effectif total, la fréquence de la valeur $2$ et la moyenne.\par\smallskip\renewcommand{\arraystretch}{1.15}\begin{tabular}{|c|c|c|c|c|}\hline
Valeur & 1 & 2 & 3 & 4 \\ \hline
Effectif & 5 & 8 & 4 & 3 \\ \hline
\end{tabular}\par }{Total $20$\,; fréquence $\dfrac{8}{20}=40\,\%$\,; moyenne $\dfrac{5+16+12+12}{20}=2{,}25$.}{}
\qitem{Fil rouge · Flash}{REP}{Donner les coordonnées du symétrique de $A(3;-2)$ par rapport à l'axe des ordonnées.}{$(-3;-2)$.}{Piège~: changer le signe de l'ordonnée (symétrie par rapport à l'axe des abscisses).}
\end{fiche}

\fichehead{jeu.~26~nov.~2026}{Intervalles de $\R$~: notations et représentation}
\begin{fiche}
\qitem{Échauffement · Flash}{CMP}{Comparer $-3$ et $-5$ puis $-\dfrac12$ et $-0{,}6$.}{$-3>-5$\,; $-0{,}5>-0{,}6$.}{Pour les négatifs, le plus petit en valeur absolue est le plus grand.}
\qitem{Cœur · Méthode}{NUM}{Écrire sous forme d'intervalle~: $-2\leqslant x<5$\,; $x>3$\,; $x\leqslant0$.}{$[-2;5[$\,; $]3;+\infty[$\,; $]-\infty;0]$.}{Le crochet est tourné vers l'intérieur si la borne est incluse\,; $\infty$ est toujours exclu.}
\qitem{Cœur + · QCM}{NUM}{L'intersection $]-\infty;2[\,\cap\,[-1;4]$ est~:\choix{$]2;4]$}{$]-\infty;4]$}{$[-1;2[$}{$[-1;2]$}}{C.\ $[-1;2[$}{Intersection~: valeurs vérifiant les deux conditions $x<2$ et $-1\leqslant x\leqslant4$.}
\qitem{Retour · Tableau}{LEC}{Durées de trajet (min).\par\smallskip\renewcommand{\arraystretch}{1.15}\begin{tabular}{|c|c|c|c|}\hline
Classe & $[0;10[$ & $[10;20[$ & $[20;30[$ \\ \hline
Effectif & 4 & 9 & 7 \\ \hline
\end{tabular}\par Quel est l'effectif total, la classe modale et la fréquence de $[10;20[$\,?}{$20$\,; $[10;20[$\,; $\dfrac{9}{20}=45\,\%$.}{}
\qitem{Fil rouge · Problème}{DIV}{Un bus passe toutes les $15$ min, un autre toutes les $20$ min. Ils se croisent à $8$ h. À quelle heure se recroisent-ils\,?}{$\ppcm(15\,;\,20)=60$~: à $9$ h.}{}
\end{fiche}

\fichehead{lun.~30~nov.~2026}{Valeur absolue et distance entre deux réels}
\begin{fiche}
\qitem{Échauffement · Flash}{REP}{Quelle est la distance entre $-3$ et $5$ sur la droite graduée\,?}{$|5-(-3)|=8$.}{}
\qitem{Cœur · Flash}{CAL}{Calculer $|{-}7|$, $|3-8|$ et $|\pi-4|$.}{$7$\,; $5$\,; $|\pi-4|=4-\pi$ car $\pi<4$.}{}
\qitem{Cœur + · Méthode}{EQ}{Résoudre $|x-3|=2$.}{Les réels à distance $2$ de $3$~: $x=1$ ou $x=5$.}{Interprétation géométrique plutôt que cas par cas.}
\qitem{Retour · Formule}{IMG}{$f(x)=x^2-3$. Donner l'image de $-2$ et les antécédents de $6$.}{$f(-2)=1$\,; $x^2-3=6\Leftrightarrow x^2=9$~: antécédents $-3$ et $3$.}{}
\qitem{Fil rouge · Formule → tableau}{REG}{$f(x)=-x+3$. Compléter~: $f(-1)$, $f(0)$, $f(2)$.}{$4\ ;\ 3\ ;\ 1$.}{}
\end{fiche}

\fichehead{mar.~1\textsuperscript{er}~déc.~2026}{Inéquation $|x-a|\leqslant r$ et intervalle associé}
\begin{fiche}
\qitem{Échauffement · Flash}{EQ}{Résoudre $2x-3\leqslant5$.}{$2x\leqslant8$, $x\leqslant4$~: $S=]-\infty;4]$.}{}
\qitem{Cœur · Méthode}{EQ}{Résoudre $|x-2|\leqslant3$ et donner $S$ sous forme d'intervalle.}{Les réels à distance au plus $3$ de $2$~: $S=[-1;5]$.}{Contrôle~: $|{-}1-2|=3$ et $|5-2|=3$.}
\qitem{Cœur + · Traduire}{NUM}{Traduire par une inégalité avec valeur absolue~: $x\in[-1;5]$ puis $x\in]4;10[$.}{$|x-2|\leqslant3$ (centre $2$, rayon $3$)\,; $|x-7|<3$ (centre $7$, rayon $3$).}{}
\qitem{Retour · Problème}{CMP}{Lot A~: $6$ yaourts pour $4{,}20$ €. Lot B~: $10$ yaourts pour $6{,}80$ €. Lequel est le plus avantageux\,?}{A~: $0{,}70$ € l'unité\,; B~: $0{,}68$ €~: B.}{}
\qitem{Fil rouge · Flash}{GEO}{Un disque a pour rayon $5$ cm. Donner son aire et son périmètre (valeurs exactes).}{Aire $=\pi\times5^2=25\pi$ cm$^2$\,; périmètre $=2\pi\times5=10\pi$ cm.}{}
\end{fiche}
